\documentclass[11pt,a4paper]{book}

\usepackage[table,dvipsnames,svgnames]{xcolor}
\usepackage{listingsutf8}
%\usepackage[utf8]{inputenc}% gestion des accents (source)
\usepackage[T1]{fontenc}% gestion des accents (PDF)
\usepackage{fontspec}
\usepackage[french]{babel}% gestion du français
\usepackage{textcomp}% caractères additionnels
\usepackage{mathtools,amssymb,amsthm}% packages de l'AMS + mathtools
\usepackage{amsmath}
\usepackage{yhmath}
\usepackage{mathpazo}
\usepackage[scaled=0.95]{helvet}
\usepackage{courier}
%\usepackage{fourier}
%\usepackage[scaled=0.95]{helvet}
%\usepackage{courier}
\usepackage{bm}
\usepackage{gensymb}
\usepackage{graphicx}

\usepackage{enumitem}
\usepackage{pifont}

\usepackage{pstricks}
\usepackage{pstricks-add}
\usepackage{pst-eucl}
\usepackage{pgfplots}
\let\clipbox\relax
%\usepackage{pst-plot,pst-text,pst-tree,pstricks-add}
%\usepackage{pst-node} % required package
\usepackage{pst-all}
\usepackage[most,skins,breakable]{tcolorbox}

\usepackage{euscript}





% ############### Nouvelles commandes ###############
% Coloration des labels des points dans pst-eucl
\makeatletter
\define@key[psset]{pst-eucl}{LabelColor}{\def\psk@LabelColor{#1}}
\psset{LabelColor=black}
\def\Pst@WhichLabel#1{{\color{\psk@LabelColor}%
		\ifx\psk@PointName\@default#1\else\psk@PointName\fi}}
\makeatother

% renomme le symbole #2 en #1#2 et invalide le symbole initial #2:
\def\renamesymbol#1#2{
	\expandafter\let\expandafter\newsym\expandafter=\csname#2\endcsname
	\expandafter\global\expandafter\let\csname#1#2\endcsname=\newsym
	\expandafter\global\expandafter\let\csname#2\endcsname=\origsym
}



\usepackage{ProfCollege}
\renamesymbol{Casio}{Calculatrice}
\usepackage[nonshellescape]{ProfLycee}



%\usepackage[left=1.5cm,right=1.5cm,top=2.5cm,bottom=2.5cm]{geometry}
\usepackage[left=13mm,right=13mm,top=18mm,bottom=18mm,footskip=30pt]{geometry}
\setlength{\headheight}{15pt}

\usepackage{titlesec}
\usepackage{pas-math}
\usepackage{longtable}
\usepackage{tikz,tkz-tab}
\usepackage{amsmath}
\usepackage[right]{eurosym}
\usepackage{sectsty}
\usepackage{tipfr}
\usepackage{setspace}

%pour faire les touches de la calculatrices
\usepackage{menukeys}
\renewmenumacro{\menu}{roundedmenus} % default: menus
\renewmenumacro{\directory}{pathswithfolder} % default: paths
\renewmenumacro{\keys}{shadowedroundedkeys} % default: roundedkeys


\usepackage{multicol}
\usepackage{multirow}
\usepackage{systeme}

%% Format des titres de chapitres, paragraphes, ...
\titleformat{\chapter}[frame]
{\color{red!85!white}\normalsize}%
{\filright\sffamily\Large%
	\enspace Chapitre \thechapter\enspace}%
{8pt}
{\color{red!85!white}\rule{0pt}{30pt}\sffamily\Huge\bfseries\filcenter}
%%%%%%%%%%%%%%%%%%%%
\titlespacing{\chapter}{0pc}{0pc}{2pc}
\titlespacing{\subsection}{2pc}{}{}

\titleformat{\subsection}[block]{\hspace{1em}}{\thesubsection}{1em .}{}
\titleformat{\subsubsection}[block]{\hspace{2em}}{\thesubsection}{1em )}{}

\chapterfont{\color{red}{}}
\sectionfont{\color{green!55!black}{}}
\subsectionfont{\color{cyan!65!black}{}}
\subsubsectionfont{\color{violet}{}}
\renewcommand{\thesection}{\color{green!55!black}\Roman{section}.}
\renewcommand{\thesubsection}{\color{cyan!65!black}\arabic{subsection})}
\renewcommand{\thesubsubsection}{\color{violet}\alph{subsubsection})}
\setcounter{secnumdepth}{3}

\titleformat{\subsection}[block]{\hspace{1em}}{\thesubsection}{1em .}{}
\titleformat{\subsubsection}[block]{\hspace{2em}}{\thesubsection}{1em )}{}

\usepackage[column=A]{cellspace}
\setlength{\cellspacebottomlimit}{4pt}
\setlength{\cellspacetoplimit}{4pt}
\setlength{\parindent}{2pt}
\setlength{\tabcolsep}{0.5cm}


%entête et pieds de pages
\usepackage{fancyhdr}
\pagestyle{fancy}
\fancyhf{}
\usepackage{lastpage}
\renewcommand{\chaptermark}[1]{\markboth {Chap. \thechapter~ $-$ ~#1}{}}
\renewcommand\headrulewidth{1pt}
\fancyhead[LO,RE]{1ère ST2S}
\fancyhead[RO,LE]{\leftmark}
\renewcommand\footrulewidth{1pt}
\fancyfoot[LO,RE]{Année 2026-2027}
\fancyfoot[RO,LE]{\textbf{Page \thepage/\pageref{LastPage}}}

%Mes boites pour les définitions, propriétés, remarques, ...

%Mes boites pour les définitions, propriétés, remarques, ...

\newcommand\titrebox[1]{\textbf{#1}}
\newtcolorbox{definition}{breakable,colback=red!10!white,
	colframe=red!85!white, fonttitle=\bfseries,enhanced,
	attach boxed title to top left={xshift=5mm, yshift=-2mm},boxed title style={size=small,colback=red!10!white,colframe=red!85!white, arc=0.5mm},
	title=\textcolor{red!85!white}{Définition} , arc=3mm}

\newtcolorbox{definitionTitre}[1]{breakable,colback=red!10!white,
	colframe=red!85!white, fonttitle=\bfseries,enhanced,
	attach boxed title to top left={xshift=5mm, yshift=-2mm},boxed title style={size=small,colback=red!10!white,colframe=red!85!white, arc=0.5mm},
	title=\textcolor{red!85!white}{Définition - #1} , arc=3mm}

\newtcolorbox{definitionPerso}[1]{breakable,colback=red!10!white,
	colframe=red!85!white, fonttitle=\bfseries,enhanced,
	attach boxed title to top left={xshift=5mm, yshift=-2mm},boxed title style={size=small,colback=red!10!white,colframe=red!85!white, arc=0.5mm},
	title=\textcolor{red!85!white}{#1} , arc=3mm}

\newtcolorbox{regle}{breakable,colback=red!10!white,
	colframe=red!85!white, fonttitle=\bfseries,enhanced,
	attach boxed title to top left={xshift=5mm, yshift=-2mm},boxed title style={size=small,colback=red!10!white,colframe=red!85!white, arc=0.5mm},
	title=\textcolor{red!85!white}{Règle} , arc=3mm}

\newtcolorbox{regles}{breakable,colback=red!10!white,
	colframe=red!85!white, fonttitle=\bfseries,enhanced,
	attach boxed title to top left={xshift=5mm, yshift=-2mm},boxed title style={size=small,colback=red!10!white,colframe=red!85!white, arc=0.5mm},
	title=\textcolor{red!85!white}{Règles} , arc=3mm}

\newtcolorbox{definitions}{breakable,colback=red!10!white,
	colframe=red!85!white, fonttitle=\bfseries,enhanced,
	attach boxed title to top left={xshift=5mm, yshift=-2mm},boxed title style={size=small,colback=red!10!white,colframe=red!85!white, arc=0.5mm},
	title=\textcolor{red!85!white}{Définitions} , arc=3mm}

\newtcolorbox{propriete}{breakable,colback=red!10!white,
	colframe=red!85!white, fonttitle=\bfseries,enhanced,
	attach boxed title to top left={xshift=5mm, yshift=-2mm},boxed title style={size=small,colback=red!10!white,colframe=red!85!white, arc=0.5mm},
	title=\textcolor{red!85!white}{Propriété} , arc=3mm}

\newtcolorbox{proprieteTitre}[1]{breakable,colback=red!10!white,
	colframe=red!85!white, fonttitle=\bfseries,enhanced,
	attach boxed title to top left={xshift=5mm, yshift=-2mm},boxed title style={size=small,colback=red!10!white,colframe=red!85!white, arc=0.5mm},
	title=\textcolor{red!85!white}{Propriété - #1} , arc=3mm}

\newtcolorbox{proprietes}{breakable,colback=red!10!white,
	colframe=red!85!white, fonttitle=\bfseries,enhanced,
	attach boxed title to top left={xshift=5mm, yshift=-2mm},boxed title style={size=small,colback=red!10!white,colframe=red!85!white, arc=0.5mm},
	title=\textcolor{red!85!white}{Propriétés} , arc=3mm}

\newtcolorbox{theoreme}{breakable,colback=red!10!white,
	colframe=red!85!white, fonttitle=\bfseries,enhanced,
	attach boxed title to top left={xshift=5mm, yshift=-2mm},boxed title style={size=small,colback=red!10!white,colframe=red!85!white, arc=0.5mm},
	title=\textcolor{red!85!white}{Théorème} , arc=3mm}

\newtcolorbox{theoremeTitre}[1]{breakable,colback=red!10!white,
	colframe=red!85!white, fonttitle=\bfseries,enhanced,
	attach boxed title to top left={xshift=5mm, yshift=-2mm},boxed title style={size=small,colback=red!10!white,colframe=red!85!white, arc=0.5mm},
	title=\textcolor{red!85!white}{Théorème - #1} , arc=3mm}

\newtcolorbox{theoremes}{breakable,colback=red!10!white,
	colframe=red!85!white, fonttitle=\bfseries,enhanced,
	attach boxed title to top left={xshift=5mm, yshift=-2mm},boxed title style={size=small,colback=red!10!white,colframe=red!85!white, arc=0.5mm},
	title=\textcolor{red!85!white}{Théorèmes} , arc=3mm}

\newtcolorbox{bexemple}{
	enhanced,
	breakable,
	fonttitle=\bfseries,
	frame empty,
	nobeforeafter,
	borderline west={0.5mm}{0mm}{cyan!75!black},
	title=\textcolor{cyan!75!black}{\faHubspot\:\:Exemple},
	%title after break=Exemple (suite),
	colback=white,
	extras={enhanced, frame empty},
	detach title,before upper={\tcbtitle\ :\quad },
}

\newtcolorbox{bexemples}{
	enhanced,
	breakable,
	fonttitle=\bfseries,
	frame empty,
	nobeforeafter,
	borderline west={0.5mm}{0mm}{cyan!75!black},
	title=\textcolor{cyan!75!black}{\faHubspot\:\:Exemples},
	%title after break=Exemples (suite),
	colback=white,
	extras={enhanced, frame empty},
	detach title,before upper={\tcbtitle\ :\quad },
}

\newtcolorbox{bexemplet}[1]{
	enhanced,
	breakable,
	fonttitle=\bfseries,
	frame empty,
	borderline west={0.5mm}{0mm}{cyan!75!black},
	title=\textcolor{cyan!75!black}{\faHubspot\:\:Exemple - #1},
	colback=white,
	extras={enhanced, frame empty},
	detach title,before upper={\tcbtitle\ :\quad },
}

\newtcolorbox{brmq}{
	enhanced,
	breakable,
	fonttitle=\bfseries,
	frame empty,
	nobeforeafter,
	borderline west={0.5mm}{0mm}{orange!95!white},
	title=\textcolor{orange!85!white}{\faHandPointRight[regular]\:\:Remarque} ,
	colback=white,
	extras={enhanced, frame empty},
	detach title,before upper={\tcbtitle\ \quad },
}

\newtcolorbox{brmqs}{
	enhanced,
	breakable,
	fonttitle=\bfseries,
	frame empty,
	nobeforeafter,
	borderline west={0.5mm}{0mm}{orange!85!white},
	title=\textcolor{orange!85!white}{\faHandPointRight[regular]\:\:Remarques},
	colback=white,
	detach title,before upper={\tcbtitle\ \quad },
	extras={enhanced, frame empty},
}

\newtcolorbox{brmqt}[1]{
	enhanced,
	breakable,
	fonttitle=\bfseries,
	frame empty,
	nobeforeafter,
	borderline west={0.5mm}{0mm}{orange!85!white},
	title=\textcolor{orange!85!white}{\faHandPointRight[regular]\:\:Remarque - #1},
	colback=white,
	extras={enhanced, frame empty},
	detach title,before upper={\tcbtitle\ \quad },
}

\newtcolorbox{consequences}[1]{
	enhanced,
	breakable,
	fonttitle=\bfseries,
	frame empty,
	nobeforeafter,
	borderline west={0.5mm}{0mm}{ProcessBlue},
	title=\textcolor{ProcessBlue}{#1},
	colback=white,
	extras={enhanced, frame empty},
	detach title,before upper={\tcbtitle\ :\quad },
}

\newtcolorbox{bmethode}{
	enhanced,
	breakable,
	fonttitle=\bfseries,
	frame empty,
	nobeforeafter,
	borderline west={0.5mm}{0mm}{blue!75!black},
	title=\textcolor{blue!75!black}{Méthode} ,
	colback=white,
	extras={enhanced, frame empty},
	attach boxed title to top left = {xshift = -3mm},
	boxed title style={size=small,colback=white,colframe=blue!75!black, arc=3mm,
	}
}

\newtcolorbox{bmethodes}{
	enhanced,
	breakable,
	fonttitle=\bfseries,
	frame empty,
	nobeforeafter,
	borderline west={0.5mm}{0mm}{blue!75!black},
	title=\textcolor{blue!75!black}{Méthodes} ,
	colback=white,
	extras={enhanced, frame empty},
	attach boxed title to top left = {xshift = -3mm},
	boxed title style={size=small,colback=blue!75!black,colframe=blue!75!black, arc=3mm,
	}
	boxed title after break style={size=small,colback=white,colframe=blue!75!black, arc=3mm,
	}
}

\newtcolorbox{bmethodet}[1]{
	enhanced,
	breakable,
	fonttitle=\bfseries,
	frame empty,
	nobeforeafter,
	borderline west={0.5mm}{0mm}{blue!75!black},
	title=\textcolor{blue!75!black}{\faTasks\:\:Méthode - #1} ,
	colback=white,
	extras={enhanced, frame empty},
	attach boxed title to top left = {xshift = -3mm},
	boxed title style={size=small,colback=white,colframe=blue!75!black, arc=3mm,
	}
}

\newtcolorbox{bdemo}{
	enhanced,
	breakable,
	fonttitle=\bfseries,
	frame empty,
	nobeforeafter,
	borderline west={0.5mm}{0mm}{violet!65!white},
	title={\faEdit[regular]\:\:Démonstration},
	%title after break=Démonstration (suite),
	colback=white,
	extras={enhanced, frame empty},
	attach boxed title to top left = {xshift = -3mm},
	boxed title style={size=small,colback=violet!65!white,colframe=violet!65!white, arc=3mm,
	}
}

\newtcolorbox{bdemos}{
	enhanced,
	breakable,
	fonttitle=\bfseries,
	frame empty,
	nobeforeafter,
	borderline west={0.5mm}{0mm}{violet!65!white},
	title={\faEdit[regular]\:\:Démonstrations},
	%title after break=Démonstration (suite),
	colback=white,
	extras={enhanced, frame empty},
	attach boxed title to top left = {xshift = -3mm},
	boxed title style={size=small,colback=violet!65!white,colframe=violet!65!white, arc=3mm,
	}
}

\newtcolorbox{bdemot}[1]{
	enhanced,
	breakable,
	fonttitle=\bfseries,
	frame empty,
	nobeforeafter,
	borderline west={0.5mm}{0mm}{violet!65!white},
	title={\faEdit[regular]\:\:Démonstration - #1},
	%title after break=Démonstration (suite),
	colback=white,
	extras={enhanced, frame empty},
	attach boxed title to top left = {xshift = -3mm},
	boxed title style={size=small,colback=violet!65!white,colframe=violet!65!white, arc=3mm,
	}
}

\newtcolorbox{boxPerso}[2]{
	enhanced,
	breakable,
	fonttitle=\bfseries,
	frame empty,
	nobeforeafter,
	borderline west={0.5mm}{0mm}{#2},
	title=#1,
	%title after break=Exemple (suite),
	colback=white,
	extras={enhanced, frame empty},
	detach title,before upper={\tcbtitle\ :\quad },
}


\newcommand*{\Coord}[3]{% 
	\ensuremath{\vv{#1}\, 
		\begin{pmatrix} 
			#2\\ 
			#3 
\end{pmatrix}}}

\newcommand*{\CoordSV}[2]{% 
	\ensuremath{\begin{pmatrix} 
			#1\\ 
			#2 
\end{pmatrix}}}

\newcommand*{\Coordp}[3]{% 
	\ensuremath{\vv{#1}\, 
		\left\lvert 
		\begin{matrix} 
			#2\\ 
			#3 
		\end{matrix}  
		\right.% Ne pas oublier le délimiteur invisible. 
}}


%personnalisation des listes non numérotées pour les différentes boites 
\newlist{itemEx}{itemize}{2}
\setlist[itemEx,1]{label=$\bullet$, font=\color{cyan!75!black}}
\setlist[itemEx,2]{label=$\bullet$, font=\color{cyan!75!black}}

\newlist{enumEx}{enumerate}{2}
\setlist[enumEx,1]{label=\arabic*), font=\bf\color{cyan!75!black}}
\setlist[enumEx,2]{label=\alph*), font=\bf\color{cyan!75!black}}

\newlist{itemMet}{itemize}{2}
\setlist[itemMet,1]{label=$\bullet$, font=\color{blue!75!black}}
\setlist[itemMet,2]{label=$\bullet$, font=\color{blue!75!black}}

\newlist{enumMet}{enumerate}{2}
\setlist[enumMet,1]{label=\arabic*), font=\bf\color{blue!75!black}}
\setlist[enumMet,2]{label=\alph*), font=\bf\color{blue!75!black}}

\newlist{itemDef}{itemize}{2}
\setlist[itemDef,1]{label=$\bullet$, font=\color{ForestGreen}}
\setlist[itemDef,2]{label=$\bullet$, font=\color{ForestGreen}}

\newlist{enumDef}{enumerate}{2}
\setlist[enumDef,1]{label=\arabic*), font=\bf\color{ForestGreen}}
\setlist[enumDef,2]{label=\alph*), font=\bf\color{ForestGreen}}

\newlist{itemProp}{itemize}{2}
\setlist[itemProp,1]{label=$\bullet$, font=\color{red!85!white}}
\setlist[itemProp,2]{label=$\bullet$, font=\color{red!85!white}}

\newlist{enumProp}{enumerate}{2}
\setlist[enumProp,1]{label=\arabic*), font=\bf\color{red!85!white}}
\setlist[enumProp,2]{label=\alph*), font=\bf\color{red!85!white}}

\newlist{itemDem}{itemize}{2}
\setlist[itemDem,1]{label=$\bullet$, font=\color{violet!85!white}}
\setlist[itemDem,2]{label=$\bullet$, font=\color{violet!85!white}}

\newlist{enumDem}{enumerate}{2}
\setlist[enumDem,1]{label=\arabic*), font=\bf\color{violet!85!white}}
\setlist[enumDem,2]{label=\alph*), font=\bf\color{violet!85!white}}

\newlist{itemRem}{itemize}{2}
\setlist[itemRem,1]{label=\ding{43}, font=\color{orange!85!white}}
\setlist[itemRem,2]{label=\ding{43}, font=\color{orange!85!white}}

\newlist{enumRem}{enumerate}{2}
\setlist[enumRem,1]{label=\arabic*), font=\bf\color{orange!85!white}}
\setlist[enumRem,2]{label=\alph*), font=\bf\color{orange!85!white}}

% nouvelle commande pour mettre entre guillemets une expression
\newcommand{\guillemets}[1]{\og \emph{#1} \fg}


\newcommand{\corr}[1]{\fcolorbox{red}{yellow!30}{#1}}

\newtcolorbox{corrBox}{enhanced,breakable,boxrule=1pt,nobeforeafter,boxsep=0pt,left=\fboxsep,right=\fboxsep,colback=yellow!30,colframe=red!75!black}

% commandes pour mettre en couleur (et en gras) dans une expression mathématique.
\newcommand{\mVert}[1]{{\color{green!55!black}\bm{#1}}}
\newcommand{\mVertWB}[1]{{\color{green!55!black}#1}}
\newcommand{\mRouge}[1]{{\color{red!85!white}\bm{#1}}}
\newcommand{\mRougeWB}[1]{{\color{red!85!white}#1}}
\newcommand{\mBleu}[1]{{\color{blue!75!black}\bm{#1}}}
\newcommand{\mBleuWB}[1]{{\color{blue!75!black}#1}}
\newcommand{\mOrange}[1]{{\color{orange!75!white}\bm{#1}}}
\newcommand{\mOrangeWB}[1]{{\color{orange!75!white}#1}}
\newcommand{\mViolet}[1]{{\color{violet!85!white}\bm{#1}}}
\newcommand{\mVioletWB}[1]{{\color{violet!85!white}#1}}

%commandes pour les flèches de distributivité.
\newcommand{\tikzmarkBB}[2][]{ \tikz[remember picture, overlay, baseline=-0.5ex]\node[#1, left](#2){}; }
\newcommand{\flechePropV}[4][]{ \tikz[remember picture, overlay]\draw (#2)edge[bend left=90,min distance=1em, -stealth]node[midway, #1]{#4}(#3); }
\newcommand{\flechePropH}[4][]{ \tikz[remember picture, overlay]\draw (#2)edge[bend left=90, min distance=1em, -stealth]node[midway, #1]{#4}(#3); }
\newcommand{\flechePropHbas}[4][]{ \tikz[remember picture, overlay]\draw (#2)edge[bend left=-90, min distance=1em, -stealth]node[midway, #1]{#4}(#3); }


\lstdefinestyle{algo}{%
	inputencoding=utf8,
	mathescape,
	basicstyle=\footnotesize\ttfamily,
	numberstyle=\scriptsize\ttfamily\color{gray},
	showstringspaces=false,
	tabsize=2,
	numbers=left,
	framexleftmargin=8mm,
	xleftmargin=8mm,
	keepspaces=true,
	frame=lines,
	framerule=1pt,
	rulecolor=\color{green!50!black},
	backgroundcolor=\color{green!5},
	keywords = {Pour, allant, Si, alors, Sinon, Fin},    
	morecomment=[l][commentstyle]{//},
	commentstyle=\color{gray},
	classoffset=0,
	keywordstyle=\color{blue!75},
	classoffset=1,
	morekeywords = {Afficher},
	keywordstyle=\color{green!50!black},
	classoffset=0,
	morestring=[b]",
	stringstyle=\color{red!75},
	literate=
	{é}{{\'e}}{1}%
	{è}{{\`e}}{1}%
	{à}{{\`a}}{1}%
	{â}{{\^a}}{1}%%%
	{ç}{{\c{c}}}{1}%
	{œ}{{\oe}}{1}%
	{ù}{{\`u}}{1}%
	{É}{{\'E}}{1}%
	{È}{{\`E}}{1}%
	{À}{{\`A}}{1}%
	{Ç}{{\c{C}}}{1}%
	{Œ}{{\OE}}{1}%
	{Ê}{{\^E}}{1}%
	{ê}{{\^e}}{1}%
	{î}{{\^i}}{1}%
	{ï}{{\"i}}{1}%%%
	{ô}{{\^o}}{1}%
	{û}{{\^u}}{1}%
	{=}{$\leftarrow$}{1}%
	{==}{$={}$}{1}}


%  --------   Macro pour les exercices et leur numérotation  

\newcounter{exo}

\setcounter{exo}{1}

\newcommand{\exercice}[1]{\vskip 0.2cm\noindent{\bf \ Exercice \theexo {  #1 }}\addtocounter{exo}{1}}

\newcounter{rit}

\setcounter{rit}{1}

\newcommand{\rituel}[1]{\vskip 0.2cm\noindent{ \bf \ \underline{Rit. \therit {  #1 }}}\addtocounter{rit}{1}}

\newcommand{\raisedrule}[2][0em]{\leaders\hbox{\rule[#1]{1pt}{#2}}\hfill}

\newcounter{numexo}

\setcounter{numexo}{1}

\newcommand{\exerciceDS}[1]{\vskip 0.6cm\noindent\rule[+2pt]{1cm}{0.5pt}\ {\large\bf \ Exercice \thenumexo {  #1 }}\raisedrule[+2pt]{0.5pt}\null\addtocounter{numexo}{1}\par\vskip 0,4cm}

\newcommand{\bareme}[1]{\tikz\node[circle,draw,inner sep=1pt,outer sep=1pt, color=red]{\textbf{#1}};}

%  --------   Pour le langage Python

% Default fixed font does not support bold face
\DeclareFixedFont{\ttb}{T1}{txtt}{bx}{n}{12} % for bold
\DeclareFixedFont{\ttm}{T1}{txtt}{m}{n}{12}  % for normal

% Custom colors
\usepackage{color}
\definecolor{deepblue}{rgb}{0,0,0.5}
\definecolor{deepred}{rgb}{0.6,0,0}
\definecolor{deepgreen}{rgb}{0,0.5,0}

\usepackage{listings}

% Python style for highlighting
\newcommand\pythonstyle{\lstset{
		language=Python,
		basicstyle=\ttm,
		otherkeywords={self},             % Add keywords here
		keywordstyle=\ttb\color{deepblue},
		emph={MyClass,__init__},          % Custom highlighting
		emphstyle=\ttb\color{deepred},    % Custom highlighting style
		stringstyle=\color{deepgreen},
		frame=tb,                         % Any extra options here
		showstringspaces=false            % 
}}


% Python environment
\lstnewenvironment{python}[1][]
{
	\pythonstyle
	\lstset{#1}
}
{}

% Python for external files
\newcommand\pythonexternal[2][]{{
		\pythonstyle
		\lstinputlisting[#1]{#2}}}

% Python for inline
\newcommand\pythoninline[1]{{\pythonstyle\lstinline!#1!}}


% \begin{document}
	
	% \section{``In-text'' listing highlighting}
	
	% \begin{python}
		% class MyClass(Yourclass):
		%     def __init__(self, my, yours):
		%         bla = '5 1 2 3 4'
		%         print bla
		% \end{python}
	
	% \section{External listing highlighting}
	
	% \pythonexternal{demo.py}
	
	% \section{Inline highlighting}
	
	% Definition \pythoninline{class MyClass} means \dots
	
	% \end{document}

\definecolor{codegreen}{rgb}{0,0.6,0}
\definecolor{codegray}{rgb}{0.5,0.5,0.5}
\definecolor{codepurple}{rgb}{0.58,0,0.82}
\definecolor{backcolour}{rgb}{0.95,0.95,0.92}

\lstdefinestyle{myStyle}{
	backgroundcolor=\color{backcolour},   
	commentstyle=\color{codegreen},
	keywordstyle=\color{magenta},
	numberstyle=\tiny\color{codegray},
	stringstyle=\color{codepurple},
	basicstyle=\ttfamily\footnotesize,
	breakatwhitespace=false,         
	breaklines=true,                 
	captionpos=b,                    
	keepspaces=true,                 
	numbers=left,                    
	numbersep=5pt,                  
	showspaces=false,                
	showstringspaces=false,
	showtabs=false,                  
	tabsize=2
}

\lstset{style=myStyle}

\newcommand{\normeB}[1]{\left\Vert \vv{#1}\right\Vert}

\newcommand\grasmaths[1]{\symbf{#1}}
\newcommand\sfbmath[1]{\symbf{\symsfup{#1}}}
\definecolor{marron}{rgb}{0.64,0.16,0.16}
\definecolor{titrebleu}{rgb}{0,0,0.6}

\usepackage{tabularx}
\usepackage{wasysym}
% ===========================tableaux de base
\newcolumntype{R}[1]{>{\raggedleft\arraybackslash }b{#1}} 
\newcolumntype{L}[1]{>{\raggedright\arraybackslash }b{#1}}
\newcolumntype{C}[1]{>{\centering\arraybackslash }b{#1}}
\newcolumntype{M}[1]{>{\centering\arraybackslash }m{#1}} %doublecentrée
\renewcommand{\tabularxcolumn}[1]{m{#1}}
\newcolumntype{Y}{>{\centering\arraybackslash }X}

%=============================boîtes de cours
\usepackage{fontawesome5}
\newcommand\fontbox{\sffamily}
%\newcommand\titrebox[1]{\textbf{#1}}
\tcbuselibrary{breakable}
\tcbset{bdefaut/.style={%
		fonttitle=\fontbox,%
		enhanced,breakable,center,boxrule=1.25pt,drop shadow=black!33!white,%
		sharp corners=uphill,arc=12pt,before skip=12pt,after skip=12pt,%
		left=5pt,leftupper=5pt,top=1mm,bottom=1mm,%
		before upper=\leavevmode
	}
}
\NewDocumentCommand\bbnewbox{ m m m m O{#4} !O{} }{%nom + titre + icone + bordure/titre/fondtitre + couleur fond + options
	\DeclareTColorBox{#1}{ O{} }{%
		bdefaut,%
		colframe=#4,coltitle=#4!25!Black,colbacktitle=#4!35,colback=#5!10,%
		title={#3\:\:\titrebox{#2##1}},%
		#6
	}%
}

\bbnewbox{bdefi}{Définition}{\faIcon[regular]{comment-dots}}{ForestGreen}
\bbnewbox{bcor}{Corollaire}{\faIcon[regular]{comment-dots}}{ForestGreen}
\bbnewbox{bdefis}{Définitions}{\faIcon[regular]{comment-dots}}{ForestGreen}
\bbnewbox{bhumour}{Humour}{\faIcon[regular]{laugh-wink}}{Yellow}
%\bbnewbox{bprop}{Propriété}{\faCog}{Salmon}
\bbnewbox{bprop}{Propriété}{\faCog}{red!85!white}
\bbnewbox{bprops}{Propriétés}{\faCog}{Salmon}
\bbnewbox{bcscq}{Conséquence}{\faArrowAltCircleRight[regular]}{Salmon}
\bbnewbox{bthm}{Théorème}{\faBullhorn}{FireBrick}
\bbnewbox{blog}{Logiciel}{\faDev}{CornflowerBlue}[white]
\bbnewbox{bhistoire}{Point histoire}{\faUniversity}{Chocolate}[white]

%\bbnewbox{brmq}{Remarque}{\faHandPointRight[regular]}{MidnightBlue}[white]
\bbnewbox{bexercice}{Exercice}{\faHubspot}{MidnightBlue}
%\bbnewbox{bexemple}{Exemple}{\faHubspot}{MediumTurquoise}
\bbnewbox{bvoc}{Vocabulaire}{\faList*[regular]}{YellowGreen}
\bbnewbox{bnota}{Notation}{\faList*}{YellowGreen}
\bbnewbox{bregle}{Règle}{\faList*}{YellowGreen}
\bbnewbox{bcadre}{Cadre}{\faInfoCircle}{Peru}
\bbnewbox{bidee}{Idée}{\faInfo}{LightBlue}
\bbnewbox{bintro}{Intro}{\faMarker}{LightSalmon}
\bbnewbox{billustr}{Illustration}{\faChartLine}{Crimson}[White]
\bbnewbox{bvideo}{}{\faIcon[regular]{file-video}}{teal}[White]
\bbnewbox{bbilan}{Bilan}{\faForward}{Purple}
%\bbnewbox{bmethode}{Méthode}{\faTasks}{Coral}
%\bbnewbox{bmethode}{Méthode}{\faTasks}{blue!75!black}
\bbnewbox{btuto}{Tuto}{\faTerminal}{Coral}[White]
\bbnewbox{bmanip}{Manipulation}{\faCodeBranch}{DarkCyan}[White]
\bbnewbox{bcalco}{Calculatrice}{\faCalculator}{MediumVioletRed}[White]
\bbnewbox{boutil}{Outil}{\faTools}{DeepSkyBlue}[White]
\bbnewbox{basavoir}{À savoir}{\faHeart}{Red}
\bbnewbox{bbloctd}{TD}{\faHourglassStart}{LightSeaGreen}[White]
\bbnewbox{bblocdm}{À vous de jouer}{\faHourglassStart}{LightSeaGreen}[White]
%\bbnewbox{brappel}{Rappel}{\faHandPointUp[regular]}{LightSalmon}
\bbnewbox{brappel}{Rappel}{\faHandPointUp[regular]}{Peru}
\bbnewbox{bresol}{Résolution}{\faPuzzlePiece}{LightSeaGreen}[White]
\bbnewbox{battention}{Attention}{\faExclamationTriangle}{Red}
\bbnewbox{bcode}{Code}{\faCode}{Crimson}[White]
\bbnewbox{bpython}{Python}{\faPython}{SeaGreen}[White]

\newcommand{\grasInDef}[1]{\textcolor{ForestGreen!55!black}{\textbf{#1}}}
\newcommand{\grasInProp}[1]{\textcolor{red!85!white}{\textbf{#1}}}
\newcommand{\grasInEx}[1]{\textcolor{MediumTurquoise!55!black}{\textbf{#1}}}
\newcommand{\grasInRem}[1]{\textcolor{orange!65!white}{\textbf{#1}}}
\newcommand{\grasInMet}[1]{\textcolor{blue!75!black}{\textbf{#1}}}
\newcommand{\grasInDem}[1]{\textcolor{violet!85!white}{\textbf{#1}}}

%\newtcolorbox{bmethodet}[1][]{%
%	bdefaut,%
%	colframe=blue!75!black,coltitle=white,colbacktitle=blue!65!black,colback=blue!75!black!15!white,
%	title={\faTasks\:\:\titrebox{{Méthode - #1}}}
%}

%outil
\newtcolorbox{boutilcalc}[1][]{%
	bdefaut,%
	colframe=DeepSkyBlue,coltitle=black,colbacktitle=DeepSkyBlue!35,colback=White,
	title=\textcolor{DeepSkyBlue!25!Black}{\faFileVideo[regular]\:\:\titrebox{Tuto#1 Vidéo}}
}

%loglogo
\newtcolorbox{bloglogo}[2][]{%
	bdefaut,%
	colframe=CornflowerBlue,coltitle=black,colbacktitle=CornflowerBlue!35,colback=white,
	title=\textcolor{CornflowerBlue!25!Black}{\faDev\:\:\titrebox{Logiciel#1}},
	overlay={\begin{tcbclipinterior}
			\node[opacity=0.33,rotate=-45] at ($(interior.south east)+(-10pt,10pt)$){\includegraphics[height=24pt]{#2}};
	\end{tcbclipinterior}}
}

%demo
%\newtcolorbox{bdemo}[1][]{%
%	bdefaut,%
%	colframe=Violet,coltitle=black,colbacktitle=Violet!35,colback=Violet!15,
%	title=\textcolor{Violet!25!Black}{\faEdit[regular]\:\:\titrebox{Démonstration#1}},
%	overlay={\begin{tcbclipinterior}%
%			\node[] at ($(interior.south east)+(-10pt,10pt)$){\large\textcolor{Violet}{\textbf{\CheckedBox}}};%
%	\end{tcbclipinterior}}
%}


%\newtcolorbox{bdemo}[1][]{%
%	bdefaut,%
%	colframe=violet!85!white,coltitle=black,colbacktitle=Violet!35,colback=Violet!15,
%	title=\textcolor{Violet!25!Black}{\faEdit[regular]\:\:\titrebox{Démonstration#1}},
%	overlay={\begin{tcbclipinterior}%
%			\node[] at ($(interior.south east)+(-10pt,10pt)$){\large\textcolor{Violet}{\textbf{\CheckedBox}}};%
%	\end{tcbclipinterior}}
%}

%demoblanc
\newtcolorbox{bdemoblanc}[1][]{%
	bdefaut,%
	colframe=Violet,coltitle=black,colbacktitle=Violet!35,colback=White,
	title=\textcolor{Violet!25!Black}{\faEdit[regular]\:\:\titrebox{Démonstration#1}},
	overlay={\begin{tcbclipinterior}%
			\node[] at ($(interior.south east)+(-10pt,10pt)$){\large\textcolor{Violet}{\textbf{\CheckedBox}}};%
	\end{tcbclipinterior}}
}

%algo
\newtcolorbox{balgo}[1][]{%
	bdefaut,%
	colframe=OrangeRed,coltitle=black,colbacktitle=OrangeRed!35,colback=White,
	title=\textcolor{OrangeRed!25!Black}{\faFileCode\:\:\titrebox{Algorithme#1}},
	overlay={\begin{tcbclipinterior}%
			\node[opacity=0.15,rotate=-15] at ($(interior.south east)+(-8pt,12pt)$){\includegraphics[height=24pt]{algorithmvec}};%
	\end{tcbclipinterior}}
}


%aide
\newtcolorbox{baide}{%
	bdefaut,%
	colframe=titrebleu!35,coltitle=titrebleu!50!Black,colbacktitle=titrebleu!15,colback=White,
	title=\textcolor{titrebleu!25!Black}{\faBacon\:\:\titrebox{Niveaux de difficulté}},
	overlay={\begin{tcbclipinterior}%
			\node[opacity=0.33,rotate=45,scale=0.33] at ($(interior.south east)+(-6pt,6pt)$){\textcolor{titrebleu!50!Black}{\footnotesize \bfseries \sffamily \textcopyright{} Dairon11}};%
	\end{tcbclipinterior}}
}

\newtcolorbox{bcalcologo}[2][]{%
	bdefaut,%
	colframe=MediumVioletRed,coltitle=black,colbacktitle=MediumVioletRed!35,colback=White,
	title=\textcolor{MediumVioletRed!25!Black}{\faCalculator\:\:\titrebox{Calculatrice#1}},
	overlay={\begin{tcbclipinterior}%
			\node[opacity=0.33,rotate=-45] at ($(interior.south east)+(-10pt,10pt)$){\includegraphics[height=48pt]{#2}};%
	\end{tcbclipinterior}}
}

\definecolor{orangescratch}{HTML}{F8A31F}
\newtcolorbox{bscratch}[1][]{%
	bdefaut,%
	colframe=orangescratch,coltitle=black,colbacktitle=orangescratch!35,colback=White,
	title=\textcolor{orange!25!Black}{\includegraphics[height=1em]{scratch}\:\:\titrebox{Scratch#1}},
	overlay={\begin{tcbclipinterior}%
			\node[opacity=0.33,rotate=45] at ($(interior.south east)+(-12pt,12pt)$){\includegraphics[height=12pt]{scratchban}};%
	\end{tcbclipinterior}}
}

%reponse
\newtcolorbox{breponse}[1][]{%
	bdefaut,%
	colframe=Thistle,coltitle=black,colbacktitle=Thistle!35,colback=White,
	before upper = \vspace{-0.05cm},
	title=\textcolor{Thistle!25!Black}{\faIcon{pen-alt}\:\:\titrebox{Réponse#1}}
}

\usepackage{dashundergaps}
\newcommand\underdash[1][black]{% <- règle la couleur du texte
	\bgroup
	\ifdim\ULdepth=\maxdimen\settodepth\ULdepth{(j}\advance\ULdepth.4pt\fi
	\markoverwith{\kern0.2em% <- règle l'espacement avant un pointillé
		\vtop{\kern0.5ex% <- règle l'altitude des pointillés
			{\color{black}% <- règle la couleur des pointillés
				\hrule width.4em% <- règle la largeur des pointillés
		}}\kern0.2em% <- règle l'espacement après un pointillé
	}\color{#1}\ULon}

\newcommand\pointille[1][1cm]{\makebox[#1]{\dotfill}}



\pgfplotsset{compat=1.18}

\begin{document}
\graphicspath{{./images/}}
% ===== Styles pgfplots (identiques à ceux du cours) =====
\pgfplotsset{
	repere/.style={
		axis lines=middle,
		axis line style={-stealth},
		grid=both,
		grid style={gray!25},
		minor tick num=0,
		tick label style={font=\scriptsize, fill=white, inner sep=1pt},
		xlabel={}, ylabel={},
		samples=150,
		clip=true,
	},
	courbe/.style={very thick, smooth},
}

\markboth{Chap. 2 $-$ Fonctions polynôme de degré 2 $-$ Exercices}{}
\thispagestyle{fancy}

\begin{center}
	{\color{red!85!white}\sffamily\Huge\bfseries Fonctions polynôme de degré 2}\\[0.2cm]
	{\color{red!85!white}\sffamily\Large Fiche d'exercices}
\end{center}

% =====================================================================
\section*{\color{green!55!black}Reconnaître une fonction polynôme de degré 2}
% =====================================================================

\exerciceDS{}

Compléter le tableau suivant.

\begin{center}
	\setlength{\extrarowheight}{3pt}
	\setlength{\tabcolsep}{0.3cm}
	\begin{tabular}{|c|c|c|c|c|c|}
		\hline
		\multirow{2}{*}{$f(x)=\ldots$} & \multicolumn{2}{c|}{Fonction polynôme de degré 2 ?} & \multirow{2}{*}{$a=\ldots$} & \multirow{2}{*}{$b=\ldots$} & \multirow{2}{*}{$c=\ldots$} \\
		\cline{2-3}
		& oui & non & & & \\
		\hline $3x^2-2x+5$ & & & & & \\
		\hline $-x^2+4$ & & & & & \\
		\hline $5x-7$ & & & & & \\
		\hline $x^3+2x^2$ & & & & & \\
		\hline $(x-2)(3x+1)$ & & & & & \\
		\hline $7-2x+0,5x^2$ & & & & & \\
		\hline $(x+1)^2-x^2$ & & & & & \\
		\hline
	\end{tabular}
\end{center}

\exerciceDS{}

Pour chacune des fonctions suivantes définies sur \R, dire s'il s'agit d'une fonction polynôme de degré 2. Si oui, donner les valeurs de $a$, $b$ et $c$.

\begin{multicols}{2}
	\begin{enumerate}[label=\textbf{\alph*)}]
		\item $f(x)=-2x^2+x-3$
		\item $g(x)=4x^2$
		\item $h(x)=1-x^2+6x$
		\item $k(x)=3x(x-4)$
		\item $m(x)=(2x-1)(2x+1)-4x^2$
		\item $p(x)=(x-3)^2$
	\end{enumerate}
\end{multicols}

% =====================================================================
\section*{\color{green!55!black}Éléments caractéristiques d'une parabole (lecture graphique)}
% =====================================================================

\exerciceDS{}

\begin{minipage}[c]{0.5\linewidth}
	On a représenté ci-contre la fonction $f$ définie sur \R par
	$$f(x)=x^2-2x-3.$$
	\begin{enumerate}[label=\textbf{\arabic*)}]
		\item Déterminer graphiquement les solutions de l'équation $f(x)=0$.
		\item Déterminer graphiquement les solutions de l'équation $f(x)=-3$.
		\item En déduire l'équation de l'axe de symétrie de la courbe.
		\item Calculer l'ordonnée du sommet $S$ de la parabole, puis donner ses coordonnées.
		\item La fonction $f$ admet-elle un maximum ou un minimum ? Préciser sa valeur.
	\end{enumerate}
\end{minipage} \hfill
\begin{minipage}[c]{0.46\linewidth}
	\begin{center}
		\begin{tikzpicture}
			\begin{axis}[repere, width=7cm, height=8cm,
				xmin=-2.6, xmax=4.6, ymin=-4.6, ymax=5.6,
				xtick={-2,...,4}, ytick={-4,...,5}]
				\addplot[courbe, green!50!black, domain=-2.6:4.6] {x^2-2*x-3};
			\end{axis}
		\end{tikzpicture}
	\end{center}
\end{minipage}

\exerciceDS{}

\begin{minipage}[c]{0.5\linewidth}
	On a représenté ci-contre la fonction $f$ définie sur \R par
	$$f(x)=-x^2+4x+1.$$
	\begin{enumerate}[label=\textbf{\arabic*)}]
		\item Déterminer graphiquement les solutions de l'équation $f(x)=1$.
		\item Tracer l'axe de symétrie de la courbe et donner son équation.
		\item Donner les coordonnées du sommet $S$ de la parabole. Justifier par un calcul.
		\item Déterminer graphiquement les solutions de l'équation $f(x)=-4$.
		\item Le coefficient $a$ est-il positif ou négatif ? Était-ce prévisible sur le graphique ?
	\end{enumerate}
\end{minipage} \hfill
\begin{minipage}[c]{0.46\linewidth}
	\begin{center}
		\begin{tikzpicture}
			\begin{axis}[repere, width=7cm, height=8cm,
				xmin=-1.6, xmax=5.6, ymin=-5.6, ymax=6.6,
				xtick={-1,...,5}, ytick={-5,...,6}]
				\addplot[courbe, green!50!black, domain=-1.6:5.6] {-x^2+4*x+1};
			\end{axis}
		\end{tikzpicture}
	\end{center}
\end{minipage}

% =====================================================================
\section*{\color{green!55!black}Associer une fonction $x\mapsto ax^2+c$ à sa courbe}
% =====================================================================

\exerciceDS{}

\begin{minipage}[c]{0.40\linewidth}
	Associer chaque fonction à sa représentation graphique. Justifier.
	$$
	\begin{aligned}
		& f(x)=2x^{2}-1 \\
		& g(x)=-x^{2}-1 \\
		& h(x)=\frac{x^{2}}{2}+2 \\
		& k(x)=-2x^{2}+2
	\end{aligned}
	$$
\end{minipage} \hfill
\begin{minipage}[c]{0.56\linewidth}
	\begin{center}
		\begin{tikzpicture}
			\begin{axis}[repere, width=8.5cm, height=8.5cm,
				xmin=-3.6, xmax=3.6, ymin=-5.6, ymax=6.6,
				xtick={-3,...,3}, ytick={-5,...,6}]
				\addplot[courbe, red!80!white, line width=1.5pt, domain=-2:2] {-2*x^2+2};
				\addplot[courbe, blue!70!white, line width=1.5pt, domain=-2:2] {2*x^2-1};
				\addplot[courbe, green!70!black, line width=1.5pt, domain=-3.6:3.6] {x^2/2+2};
				\addplot[courbe, orange, line width=1.5pt, domain=-3:3] {-x^2-1};
				\node[red!80!white, anchor=east] at (axis cs:-1.95,-5.2) {$\mathcal{C}_1$};
				\node[blue!70!white, anchor=east] at (axis cs:1.7,5) {$\mathcal{C}_2$};
				\node[green!70!black, anchor=south] at (axis cs:-3,6.5) {$\mathcal{C}_3$};
				\node[orange, anchor=west] at (axis cs:2.15,-4.6) {$\mathcal{C}_4$};
			\end{axis}
		\end{tikzpicture}
	\end{center}
\end{minipage}

\exerciceDS{}

\begin{minipage}[c]{0.40\linewidth}
	\begin{enumerate}[label=\textbf{\arabic*)}]
		\item Associer chaque fonction à sa représentation graphique. Justifier.
		$$
		\begin{aligned}
			& f(x)=x^{2}-4 \\
			& g(x)=-\frac{x^{2}}{2}+3 \\
			& h(x)=3x^{2} \\
			& k(x)=-3x^{2}
		\end{aligned}
		$$
		\item Pour chacune de ces fonctions, donner l'axe de symétrie et les coordonnées du sommet de sa courbe.
	\end{enumerate}
\end{minipage} \hfill
\begin{minipage}[c]{0.56\linewidth}
	\begin{center}
		\begin{tikzpicture}
			\begin{axis}[repere, width=8.5cm, height=8.5cm,
				xmin=-4.6, xmax=4.6, ymin=-5.6, ymax=5.6,
				xtick={-4,...,4}, ytick={-5,...,5}]
				\addplot[courbe, red!80!white, line width=1.5pt, domain=-4.6:4.6] {-x^2/2+3};
				\addplot[courbe, blue!70!white, line width=1.5pt, domain=-1.5:1.5] {3*x^2};
				\addplot[courbe, green!70!black, line width=1.5pt, domain=-3.2:3.2] {x^2-4};
				\addplot[courbe, orange, line width=1.5pt, domain=-1.5:1.5] {-3*x^2};
				\node[red!80!white, anchor=south east] at (axis cs:-3.35,-2.3) {$\mathcal{C}_1$};
				\node[blue!70!white, anchor=east] at (axis cs:-1.3,5) {$\mathcal{C}_2$};
				\node[green!70!black, anchor=west] at (axis cs:3,4.5) {$\mathcal{C}_3$};
				\node[orange, anchor=east] at (axis cs:-1.3,-5) {$\mathcal{C}_4$};
			\end{axis}
		\end{tikzpicture}
	\end{center}
\end{minipage}

% =====================================================================
\section*{\color{green!55!black}Tableau de variations et sommet}
% =====================================================================

\exerciceDS{}

On donne ci-dessous le tableau de variations d'une fonction polynôme du second degré $f$ définie sur \R par $f(x)=ax^2+bx+c$.

\begin{center}
	\begin{tikzpicture}
		\tikzset{node style/.append style = {fill=none}}
		\tkzTabInit[lgt = 1.6, espcl = 2.4, deltacl = 0.5]{$x$ / 1 , $f(x)$ / 2}{$-\infty$, $-2$, $+\infty$}
		\tkzTabVar{-/, +/ $5$, -/ }
	\end{tikzpicture}
\end{center}

\begin{enumerate}[label=\textbf{\arabic*)}]
	\item Le coefficient $a$ est-il positif ou négatif ? Justifier.
	\item Donner les coordonnées du sommet $S$ de la parabole représentant $f$.
	\item Donner l'équation de l'axe de symétrie de la parabole.
	\item La fonction $f$ admet-elle un maximum ou un minimum sur \R ? Préciser sa valeur et pour quelle valeur de $x$ il est atteint.
	\item Comparer, sans calcul, $f(-5)$ et $f(-3)$.
	\item On sait que $f(0)=1$. En utilisant la symétrie de la parabole, donner la valeur de $f(-4)$.
\end{enumerate}

\exerciceDS{}

On considère la fonction $f$ définie sur \R par $f(x)=2x^2-8x+3$.

\begin{enumerate}[label=\textbf{\arabic*)}]
	\item Donner les valeurs de $a$, $b$ et $c$. En déduire le sens de variation de $f$ (d'abord croissante puis décroissante, ou l'inverse).
	\item Calculer $f(1)$ et $f(3)$.
	\item En déduire l'équation de l'axe de symétrie de la parabole représentant $f$.
	\item Calculer les coordonnées du sommet $S$ de la parabole.
	\item Dresser le tableau de variations de la fonction $f$.
\end{enumerate}

% =====================================================================
\section*{\color{green!55!black}Forme factorisée et forme développée}
% =====================================================================

\exerciceDS{}

Pour chacune des fonctions suivantes définies sur \R :
\begin{itemize}
	\item développer l'expression et donner les valeurs de $a$, $b$ et $c$ ;
	\item donner les racines de la fonction polynôme.
\end{itemize}

\begin{multicols}{3}
	\begin{enumerate}[label=\textbf{\alph*)}]
		\item $f(x)=3(x-2)(x+4)$
		\item $g(x)=-(x+1)(x-5)$
		\item $h(x)=2(x-3)^2$
	\end{enumerate}
\end{multicols}

\exerciceDS{}

\begin{enumerate}[label=\textbf{\arabic*)}]
	\item Développer les expressions suivantes :
	\begin{multicols}{2}
		\begin{enumerate}[label=\textbf{\alph*)}]
			\item $f(x)=-2(x+3)(x-1)$
			\item $g(x)=\dfrac{1}{2}(x-4)(x+2)$
		\end{enumerate}
	\end{multicols}
	\item Donner les solutions des équations $f(x)=0$ et $g(x)=0$.
	\item Écrire sous forme factorisée la fonction polynôme $p$ de degré 2 dont les racines sont $2$ et $-5$ et dont le coefficient $a$ est égal à $3$. Développer ensuite l'expression obtenue.
\end{enumerate}

% =====================================================================
\section*{\color{green!55!black}Associer une forme factorisée à sa courbe}
% =====================================================================

\exerciceDS{}

\begin{minipage}[c]{0.40\linewidth}
	Associer chaque fonction à sa représentation graphique. Justifier.
	$$
	\begin{aligned}
		& f(x)=(x-2)(x+1) \\
		& g(x)=-(x-2)(x+1) \\
		& h(x)=2(x+1)^{2}
	\end{aligned}
	$$
\end{minipage} \hfill
\begin{minipage}[c]{0.56\linewidth}
	\begin{center}
		\begin{tikzpicture}
			\begin{axis}[repere, width=8cm, height=8.5cm,
				xmin=-3.6, xmax=3.6, ymin=-6.5, ymax=8.5,
				xtick={-3,...,3}, ytick={-6,-4,...,8}]
				\addplot[courbe, red!80!white, line width=1.5pt, domain=-3:1] {2*(x+1)^2};
				\addplot[courbe, blue!70!white, line width=1.5pt, domain=-2.5:3.5] {-(x-2)*(x+1)};
				\addplot[courbe, green!70!black, line width=1.5pt, domain=-2.9:3.9] {(x-2)*(x+1)};
				\node[red!80!white, anchor=east] at (axis cs:-2.9,7) {$\mathcal{C}_1$};
				\node[blue!70!white, anchor=west] at (axis cs:-2.3,-5.5) {$\mathcal{C}_2$};
				\node[green!70!black, anchor=west] at (axis cs:3.3,6.5) {$\mathcal{C}_3$};
			\end{axis}
		\end{tikzpicture}
	\end{center}
\end{minipage}

\exerciceDS{}

\begin{minipage}[c]{0.40\linewidth}
	On a représenté trois fonctions polynômes du second degré. Parmi les quatre fonctions ci-dessous, une seule ne correspond à aucune courbe.
	$$
	\begin{aligned}
		& f(x)=\frac{1}{2}(x+2)(x-4) \\
		& g(x)=-(x+2)(x-4) \\
		& h(x)=-(x-1)^{2} \\
		& k(x)=(x-2)(x+4)
	\end{aligned}
	$$
	\begin{enumerate}[label=\textbf{\arabic*)}]
		\item Associer chaque courbe à sa fonction. Justifier.
		\item Quelle fonction est l'intruse ? Pourquoi ?
	\end{enumerate}
\end{minipage} \hfill
\begin{minipage}[c]{0.56\linewidth}
	\begin{center}
		\begin{tikzpicture}
			\begin{axis}[repere, width=8.5cm, height=9cm,
				xmin=-3.6, xmax=5.6, ymin=-7, ymax=10.5,
				xtick={-3,...,5}, ytick={-6,-4,...,10}]
				\addplot[courbe, red!80!white, line width=1.5pt, domain=-3:5] {-(x+2)*(x-4)};
				\addplot[courbe, blue!70!white, line width=1.5pt, domain=-1.5:3.5] {-(x-1)^2};
				\addplot[courbe, green!70!black, line width=1.5pt, domain=-3.6:5.6] {0.5*(x+2)*(x-4)};
				\node[red!80!white, anchor=south] at (axis cs:1,9.1) {$\mathcal{C}_1$};
				\node[blue!70!white, anchor=west] at (axis cs:3.4,-5.8) {$\mathcal{C}_2$};
				\node[green!70!black, anchor=west] at (axis cs:4.8,4.5) {$\mathcal{C}_3$};
			\end{axis}
		\end{tikzpicture}
	\end{center}
\end{minipage}

% =====================================================================
\section*{\color{green!55!black}Factoriser connaissant une racine}
% =====================================================================

\exerciceDS{}

On considère la fonction $f$ définie sur \R par $f(x)=3x^2-3x-6$.

\begin{enumerate}[label=\textbf{\arabic*)}]
	\item Vérifier que $2$ est une racine de la fonction polynôme $f$.
	\item En déduire que $f$ peut s'écrire sous la forme $f(x)=3(x-2)(x-x_2)$.
	\item En choisissant une valeur de $x$ judicieuse, déterminer $x_2$.
	\item Écrire $f$ sous forme factorisée et vérifier le résultat en développant.
\end{enumerate}

\exerciceDS{}

On considère la fonction $f$ définie sur \R par $f(x)=-2x^2+2x+12$.

\begin{enumerate}[label=\textbf{\arabic*)}]
	\item Calculer $f(-1)$, $f(1)$ et $f(3)$. Conjecturer une racine de la fonction polynôme $f$.
	\item Factoriser $f$.
	\item En déduire les solutions de l'équation $f(x)=0$.
\end{enumerate}

% =====================================================================
\section*{\color{green!55!black}Signe d'une fonction polynôme de degré 2}
% =====================================================================

\exerciceDS{}

On considère la fonction $f$ définie sur \R par $f(x)=3(x-1)(x+4)$.

\begin{enumerate}[label=\textbf{\arabic*)}]
	\item Déterminer les racines de la fonction polynôme $f$.
	\item Dresser le tableau de signes de $f(x)$ sur \R.
	\item En déduire les solutions de l'inéquation $f(x)\geq 0$, puis celles de l'inéquation $f(x)<0$.
\end{enumerate}

\exerciceDS{}

Une entreprise fabrique des objets. On admet que son bénéfice, en milliers d'euros, pour la fabrication de $x$ centaines d'objets (avec $x\in[0 ; 10]$), est donné par :
$$B(x)=-2(x-2)(x-8).$$

\begin{enumerate}[label=\textbf{\arabic*)}]
	\item Calculer $B(0)$ et $B(5)$. Interpréter ces résultats dans le contexte de l'exercice.
	\item Déterminer les racines de la fonction polynôme $B$.
	\item Dresser le tableau de signes de $B(x)$ sur l'intervalle $[0 ; 10]$.
	\item Pour quelles quantités d'objets fabriqués l'entreprise réalise-t-elle un bénéfice positif ?
	\item En utilisant la symétrie de la parabole, déterminer la quantité d'objets à fabriquer pour obtenir un bénéfice maximal. Quel est ce bénéfice maximal ?
\end{enumerate}

\end{document}
